% TOP_PROBLEM: {"id": 1100803, "title": "Questions in Geometric Group Theory — Q 8.3", "category_id": 17, "category": {"id": 17, "name": "group_theory", "display_name": "Group Theory", "description": "Problems about groups, group actions, representations, and related algebraic structures.", "slug": "group-theory", "order_index": 17, "created_at": "2026-07-31T15:26:25.671Z"}, "status": "partially_solved", "problem_number": "AMR-010-0803", "set_id": 13}
% FIRST_POSTED: 2026-10-01
% ENTRY_KIND: statement_only
% UnsolvedMath ID 1100803; source catalogue code AMR-010-0803
% !TeX program = lualatex
% Definition and sources only; this file is not an LLM solution attempt.
\documentclass[11pt,a4paper]{article}
\usepackage[margin=25mm]{geometry}
\usepackage{fontspec}
\setmainfont{lmroman10-regular.otf}[BoldFont=lmroman10-bold.otf,
 ItalicFont=lmroman10-italic.otf,BoldItalicFont=lmroman10-bolditalic.otf]
\usepackage{amsmath,amssymb,mathtools}
\usepackage{unicode-math}
\setmathfont{latinmodern-math.otf}
\AtBeginDocument{\let\setminus\smallsetminus}
\usepackage{xurl,hyperref}
\hypersetup{colorlinks=true,urlcolor=blue}
\setcounter{secnumdepth}{0}
\setlength{\parindent}{0pt}
\setlength{\parskip}{5pt}
\setlength{\emergencystretch}{3em}
\begin{document}
\section*{Questions in Geometric Group Theory — Q 8.3}\label{1100803}
{\small UnsolvedMath ID 1100803 · AMR-010-0803 · Group Theory · AMR Open Problem Lists}
\subsection{Definitions and mathematical statement}
Let $G$ be finitely presented. A simplicial tree is a connected graph with no cycles, given its unit-edge path metric. An action is nontrivial when no point of its geometric realization is fixed by every group element; subdivision allows one to remove edge inversions.

Compare two properties. (A) There is a subgroup $H\leq G$ of finite index and an action of $H$ on a simplicial tree with no global fixed point. (B) For every connected finite CW complex $X$ with $\pi_1(X)\cong G$, there is a connected covering space $\widehat X\to X$ having at least two ends.

For a connected locally finite complex, having at least two ends means that some compact subset has a complement with at least two components which are not relatively compact. Equivalently, one can use the unbounded complementary components in a proper lifted length metric. Coverings of a fixed finite complex have this standard end notion.

To what extent are (A) and (B) equivalent? Determine whether each implication holds for all finitely presented groups, and identify any necessary qualification if not. The cover in (B) may have infinitely many sheets and need not be regular, whereas the subgroup in (A) specifically has finite index. No condition on edge stabilizers of the tree action is imposed. The task concerns these two existence properties, not a claim that each multi-ended covering must itself arise from one specified splitting.
\subsection{Short English statement}
Is virtual nontrivial action on a tree equivalent to the existence of a multi-ended covering of every finite complex realizing the group?
\subsection{Sources}
\begin{itemize}
\item[S1] ULAM.AI and the UnsolvedMath Contributors, supplied problems.json, record 1100803 (AMR-010-0803), AMR Open Problem Lists. Catalogue identity and source transcription; CC BY 4.0.
\url{https://huggingface.co/datasets/ulamai/UnsolvedMath}.
\item[S2] Bestvina - Questions in Geometric Group Theory (pdf) (2004), Question 8.3 (PDF page 16). Original source indexed by AMR-010-0803.
\url{https://www.math.utah.edu/~bestvina/eprints/questions-updated.pdf}.
\end{itemize}
\end{document}
